\section{Lịch sử phát triển của Trí tuệ nhân tạo}

Trí tuệ nhân tạo (Artificial Intelligence - AI) là một trong những lĩnh vực nghiên cứu năng động và phát triển nhanh nhất trong khoa học máy tính. Mặc dù thuật ngữ "Trí tuệ nhân tạo" mới chỉ chính thức được đưa ra vào giữa thế kỷ 20, nhưng ước mơ về việc tạo ra những cỗ máy có khả năng suy nghĩ và hành động như con người đã tồn tại từ thời cổ đại qua những câu chuyện thần thoại, truyền thuyết và các tác phẩm cơ khí sơ khai. Sự phát triển của AI không phải là một đường thẳng đi lên, mà là một chuỗi các giai đoạn với những bước tiến đột phá đan xen với những kỳ vọng bị phá vỡ, thường được gọi là các "Mùa đông AI". Chương này sẽ đi sâu vào lịch sử hình thành, sự phát triển qua các thời kỳ, những cột mốc quan trọng và trạng thái hiện tại của Trí tuệ nhân tạo, đồng thời xem xét chi tiết về các cơ sở toán học và thuật toán trong những thời kỳ đầu.

Sự hình thành của AI bắt nguồn từ các lĩnh vực triết học, toán học, tâm lý học, thần kinh học và khoa học máy tính. Những câu hỏi cốt lõi như "Tâm trí hoạt động như thế nào?", "Làm thế nào một hệ thống vật lý có thể tạo ra trí thông minh?" đã thôi thúc các nhà nghiên cứu tìm kiếm cách mô phỏng các quá trình nhận thức của con người bằng các thiết bị tính toán cơ học và điện tử. Trong thế kỷ 20, sự ra đời của máy tính điện tử đã cung cấp nền tảng phần cứng cần thiết để biến những lý thuyết về trí tuệ nhân tạo thành hiện thực. 

Lịch sử AI có thể được chia thành nhiều giai đoạn chính: sự khởi đầu khái niệm với Alan Turing, Hội nghị Dartmouth huyền thoại đánh dấu sự ra đời của lĩnh vực học thuật, những thập kỷ của sự lạc quan tột độ (và sự sụp đổ sau đó), sự trỗi dậy của các hệ chuyên gia, sự trở lại của mạng nơ-ron nhân tạo nhờ thuật toán lan truyền ngược, và cuối cùng là kỷ nguyên bùng nổ của Học sâu (Deep Learning) và các mô hình ngôn ngữ lớn (Large Language Models) mà chúng ta đang chứng kiến ngày nay.

\newpage

\subsection{Sự ra đời của Trí tuệ nhân tạo (Trước 1956 và Hội nghị Dartmouth)}

Thời kỳ trước năm 1956 đánh dấu những bước chuẩn bị quan trọng về mặt lý thuyết và khái niệm cho sự ra đời của AI. Một trong những cột mốc quan trọng nhất là bài báo kinh điển "Computing Machinery and Intelligence" được xuất bản năm 1950 bởi nhà toán học, logic học và mật mã học vĩ đại người Anh, Alan Turing. Trong bài báo này, Turing đã đặt ra một câu hỏi nền tảng: "Máy móc có thể suy nghĩ không?" (Can machines think?). Để tránh những tranh cãi triết học phức tạp về định nghĩa của "suy nghĩ", Turing đã đề xuất một bài kiểm tra thực nghiệm, sau này được biết đến rộng rãi với tên gọi "Phép thử Turing" (Turing Test).

Phép thử Turing được thiết kế như một "Trò chơi bắt chước" (Imitation Game), trong đó một người thẩm vấn (interrogator) sẽ giao tiếp với một con người và một máy tính thông qua văn bản (text) ở hai phòng tách biệt. Nếu người thẩm vấn không thể phân biệt được đâu là máy tính và đâu là con người sau một khoảng thời gian giao tiếp, máy tính đó được coi là đã vượt qua phép thử và có khả năng thể hiện trí thông minh tương đương con người. Bài báo của Turing không chỉ giới thiệu Phép thử Turing mà còn thảo luận về nhiều vấn đề then chốt của AI như học máy (machine learning), thuật toán di truyền (genetic algorithms) và học củng cố (reinforcement learning).

Bên cạnh Alan Turing, những nghiên cứu đầu tiên về mạng nơ-ron nhân tạo cũng được tiến hành trong giai đoạn này. Năm 1943, Warren McCulloch (nhà thần kinh học) và Walter Pitts (nhà logic học) đã xuất bản bài báo "A Logical Calculus of the Ideas Immanent in Nervous Activity", đề xuất một mô hình toán học đơn giản đầu tiên của nơ-ron sinh học. Họ chứng minh rằng các mạng bao gồm các nơ-ron nhân tạo này có thể tính toán bất kỳ hàm logic nào. 

Về mặt toán học, mô hình McCulloch-Pitts có thể được biểu diễn như sau. Mỗi nơ-ron nhận các tín hiệu đầu vào nhị phân $x_i \in \{0, 1\}$. Trọng số của các kết nối, $w_i$, có thể là kích thích (excitatory, $w_i > 0$) hoặc ức chế (inhibitory, $w_i < 0$). Trạng thái đầu ra $y$ của nơ-ron được xác định bởi một hàm bước (step function):

\begin{equation}
y = \begin{cases} 
1 & \text{nếu } \sum_{i=1}^{n} w_i x_i \geq \theta \\
0 & \text{nếu } \sum_{i=1}^{n} w_i x_i < \theta 
\end{cases}
\end{equation}

trong đó $\theta$ là ngưỡng kích hoạt (threshold). Mặc dù đơn giản, mô hình này đã đặt nền móng lý thuyết quan trọng cho cấu trúc của các mạng nơ-ron hiện đại.

Năm 1956 được coi là năm sinh chính thức của lĩnh vực Trí tuệ nhân tạo với "Dự án Nghiên cứu Mùa hè về Trí tuệ nhân tạo tại Dartmouth" (Dartmouth Summer Research Project on Artificial Intelligence). Hội nghị này được tổ chức bởi John McCarthy, Marvin Minsky, Nathaniel Rochester và Claude Shannon. John McCarthy đã chính thức đề xuất việc sử dụng thuật ngữ "Artificial Intelligence" (Trí tuệ nhân tạo) trong đề xuất của dự án. Đề xuất của hội nghị có tuyên bố đầy tham vọng: "Nghiên cứu sẽ được tiến hành dựa trên phỏng đoán rằng mọi khía cạnh của việc học tập hoặc bất kỳ đặc điểm nào khác của trí thông minh đều có thể được mô tả chính xác đến mức một cỗ máy có thể được chế tạo để mô phỏng nó". 

Hội nghị Dartmouth quy tụ những bộ óc sáng giá nhất thời bấy giờ, trong đó có Allen Newell và Herbert A. Simon, những người đã trình bày chương trình "Logic Theorist" - chương trình máy tính đầu tiên được thiết kế để bắt chước kỹ năng giải quyết vấn đề của con người, và đã chứng minh thành công nhiều định lý trong cuốn Principia Mathematica. Sự kiện này đã thống nhất một cộng đồng nhỏ các nhà nghiên cứu xung quanh một mục tiêu chung và khởi xướng nhiều thập kỷ nghiên cứu chuyên sâu về AI.

\newpage

\subsection{Kỳ vọng ban đầu và Mùa đông AI đầu tiên (1956 - 1980)}

Những năm tiếp theo sau Hội nghị Dartmouth chứng kiến một sự lạc quan tột độ và những khoản đầu tư khổng lồ từ các cơ quan chính phủ, đặc biệt là DARPA (Defense Advanced Research Projects Agency) của Hoa Kỳ. Những người tiên phong trong lĩnh vực AI tin rằng máy móc có trí thông minh ngang ngửa con người sẽ xuất hiện trong vòng một vài thập kỷ.

Trong giai đoạn này, nhiều chương trình máy tính đáng kinh ngạc đã được phát triển. Ví dụ tiêu biểu là General Problem Solver (GPS) do Newell và Simon phát triển năm 1957. GPS được thiết kế để giải quyết các vấn đề đa dạng theo cách tư duy của con người, sử dụng phương pháp phân tích phương tiện-mục đích (means-ends analysis) để giảm thiểu sự khác biệt giữa trạng thái hiện tại và trạng thái đích. Mặc dù GPS giải quyết tốt các bài toán logic đơn giản, nó gặp khó khăn khi đối mặt với các vấn đề trong thế giới thực đòi hỏi kiến thức nền rộng lớn.

Một thành tựu nổi bật khác là chương trình ELIZA, do Joseph Weizenbaum tại MIT phát triển từ năm 1964 đến 1966. ELIZA là một trong những chương trình xử lý ngôn ngữ tự nhiên đầu tiên, được thiết kế để mô phỏng một nhà trị liệu tâm lý Rogerian. Mặc dù ELIZA chỉ sử dụng các kỹ thuật đối sánh mẫu (pattern matching) và thay thế văn bản đơn giản mà không thực sự "hiểu" ngữ nghĩa, nó đã tạo ra ảo giác mạnh mẽ về trí thông minh và sự đồng cảm, khiến nhiều người dùng tin rằng họ đang trò chuyện với một thực thể có trí óc.

Tuy nhiên, sự lạc quan quá mức đã dẫn đến những hứa hẹn phi thực tế. Khi những khó khăn thực sự của việc giải quyết các bài toán AI phức tạp lộ diện, tiến độ chậm lại đáng kể. Các hệ thống AI thời kỳ này chủ yếu dựa trên logic biểu tượng (symbolic logic) và các quy tắc tìm kiếm heuristic. Mặc dù phương pháp này có hiệu quả với các "thế giới khối" (micro-worlds) bị giới hạn, nó hoàn toàn thất bại khi đối mặt với "sự bùng nổ tổ hợp" (combinatorial explosion) trong các vấn đề thế giới thực quy mô lớn. Số lượng các trạng thái có thể xảy ra trong các bài toán thực tế tăng lên theo cấp số nhân, vượt xa khả năng tính toán của máy tính và bộ nhớ phần cứng lúc bấy giờ.

Hơn nữa, một rào cản lớn khác là "nghịch lý Moravec" (Moravec's paradox). Nghịch lý này phát biểu rằng các nhiệm vụ tư duy cấp cao (như giải toán logic, chơi cờ) đòi hỏi rất ít tính toán, trong khi các nhiệm vụ nhận thức và vận động cấp thấp (như nhận dạng hình ảnh, di chuyển trong không gian, hiểu ngôn ngữ giao tiếp hàng ngày) lại đòi hỏi những nguồn lực tính toán khổng lồ. 

Báo cáo Lighthill (Lighthill Report) năm 1973 tại Anh đã đưa ra một đánh giá rất bi quan về triển vọng của AI, chỉ trích lĩnh vực này vì không đạt được những mục tiêu vĩ đại đã hứa hẹn. Báo cáo này dẫn đến việc chính phủ Anh cắt giảm nghiêm trọng tài trợ cho các nghiên cứu AI tại hầu hết các trường đại học. Cùng thời điểm đó ở Mỹ, DARPA cũng rút bớt tài trợ cho nghiên cứu AI cơ bản do thất vọng về kết quả của dự án dịch máy tự động. Sự cắt giảm tài trợ hàng loạt và sự suy giảm niềm tin này đã mở ra "Mùa đông AI đầu tiên" (The First AI Winter), một giai đoạn kéo dài từ giữa những năm 1970 đến khoảng năm 1980, đặc trưng bởi sự đình trệ và thiếu nguồn lực trong nghiên cứu trí tuệ nhân tạo.

\newpage

\subsection{Hệ chuyên gia và Mùa đông AI thứ hai (1980 - 1993)}

Vào đầu những năm 1980, AI bắt đầu hồi sinh nhờ một hướng đi mới và thực dụng hơn: các Hệ chuyên gia (Expert Systems). Thay vì cố gắng tạo ra một trí thông minh tổng quát (Artificial General Intelligence - AGI) như trước đây, các nhà nghiên cứu nhận ra rằng họ có thể tạo ra các hệ thống hữu ích bằng cách giới hạn phạm vi ứng dụng vào các miền tri thức hẹp và cụ thể.

Một hệ chuyên gia là một hệ thống máy tính mô phỏng khả năng ra quyết định của một chuyên gia con người trong một lĩnh vực cụ thể. Hệ thống bao gồm hai thành phần chính:
\begin{itemize}
    \item \textbf{Cơ sở tri thức (Knowledge Base):} Chứa đựng kiến thức chuyên môn dưới dạng các tập luật (thường là các câu điều kiện IF-THEN) và sự kiện do các chuyên gia cung cấp.
    \item \textbf{Động cơ suy diễn (Inference Engine):} Các thuật toán logic áp dụng các luật trong cơ sở tri thức vào dữ liệu đầu vào để suy luận và đưa ra câu trả lời hoặc khuyến nghị.
\end{itemize}

Hệ chuyên gia đầu tiên thành công vang dội là R1 (sau này được đổi tên thành XCON), được phát triển bởi John McDermott tại Đại học Carnegie Mellon vào năm 1980 cho Tập đoàn Thiết bị Kỹ thuật số (Digital Equipment Corporation - DEC). XCON tự động hóa quá trình cấu hình các hệ thống máy tính VAX của DEC dựa trên yêu cầu của khách hàng, giúp DEC tiết kiệm hàng chục triệu USD mỗi năm bằng cách giảm thiểu lỗi trong các đơn đặt hàng. Một ví dụ nổi tiếng khác là MYCIN, phát triển tại Đại học Stanford, được thiết kế để chẩn đoán các bệnh nhiễm khuẩn máu và đề xuất liệu pháp kháng sinh. MYCIN có độ chính xác chẩn đoán vượt trội so với nhiều bác sĩ nội trú.

Sự thành công thương mại của các hệ chuyên gia đã kích hoạt một làn sóng đầu tư mới khổng lồ vào AI từ các tập đoàn công nghiệp. Một ngành công nghiệp hoàn toàn mới xoay quanh phần cứng và phần mềm cho AI đã ra đời, trong đó đáng chú ý nhất là thị trường máy tính Lisp (Lisp machines) – những máy tính được thiết kế đặc biệt để tối ưu hóa việc chạy ngôn ngữ lập trình Lisp phổ biến trong AI thời đó, do các công ty như Symbolics và Lisp Machines Inc. sản xuất. Chính phủ Nhật Bản cũng công bố "Dự án Máy tính Thế hệ thứ năm" (Fifth Generation Computer Systems) với khoản ngân sách hàng trăm triệu đô la, nhằm phát triển siêu máy tính có khả năng suy luận logic khổng lồ.

Tuy nhiên, thời kỳ hoàng kim của hệ chuyên gia không kéo dài. Đến cuối những năm 1980, các công ty bắt đầu nhận ra những giới hạn chí mạng của hướng tiếp cận này:
\begin{enumerate}
    \item \textbf{Khó khăn trong việc thu nhận tri thức (Knowledge Acquisition Bottleneck):} Việc trích xuất kiến thức từ các chuyên gia con người và mã hóa chúng thành các tập luật là quá trình tốn cực kỳ nhiều thời gian, công sức và tốn kém.
    \item \textbf{Tính giòn gãy (Brittleness):} Hệ chuyên gia không có "ý thức" thông thường. Khi đối mặt với một tình huống nằm ngoài các quy tắc đã được lập trình sẵn, hệ thống sẽ thất bại hoàn toàn thay vì suy thoái dần (graceful degradation) như trí tuệ con người.
    \item \textbf{Chi phí bảo trì cao:} Khi lĩnh vực chuyên môn thay đổi, cơ sở tri thức phải được cập nhật thủ công, điều này rất phức tạp đối với các hệ thống có hàng chục ngàn tập luật.
    \item \textbf{Sự sụp đổ của phần cứng chuyên dụng:} Sự phát triển thần tốc của các máy tính cá nhân (PC) và máy trạm (workstation) thông thường từ Apple và IBM đã khiến những chiếc Lisp machines đắt đỏ trở nên lỗi thời. Các máy PC thông thường có thể chạy phần mềm Lisp nhanh hơn và rẻ hơn các máy Lisp chuyên dụng.
\end{enumerate}

Sự thất vọng của giới công nghiệp khi các dự án hệ chuyên gia thất bại trong việc đáp ứng kỳ vọng, cùng với sự kết thúc tài trợ của chính phủ, đã dẫn đến sự sụp đổ của thị trường máy tính Lisp vào năm 1987 và gây ra "Mùa đông AI thứ hai". Thời kỳ này thậm chí còn tồi tệ hơn lần đầu tiên, với cụm từ "Trí tuệ nhân tạo" gần như trở thành một từ cấm kỵ trong các hồ sơ xin tài trợ nghiên cứu hoặc trong kinh doanh, buộc các nhà nghiên cứu phải ngụy trang công việc của họ dưới các tên gọi khác như "Tin học học", "Hệ thống thông minh", hoặc "Học máy".

\newpage

\subsection{Sự trỗi dậy của Thuyết kết nối và Học máy (1990s - 2010s)}

Trong khi logic biểu tượng (Symbolic AI hay GOFAI - Good Old-Fashioned AI) thống trị phần lớn thời gian trước những năm 1990, một trường phái nghiên cứu khác, được gọi là "Thuyết kết nối" (Connectionism), vẫn âm thầm phát triển. Thuyết kết nối cố gắng mô phỏng mạng lưới sinh học của não bộ thông qua các Mạng Nơ-ron Nhân tạo (Artificial Neural Networks - ANNs).

Vào năm 1958, Frank Rosenblatt đã phát minh ra Perceptron, một mô hình nơ-ron đơn giản có khả năng học hỏi qua dữ liệu để thực hiện phân loại nhị phân. Công thức toán học cốt lõi cho đầu ra $\hat{y}$ của một Perceptron cho một mẫu dữ liệu đầu vào có đặc trưng biểu diễn dưới dạng vector $x = [x_1, x_2, \dots, x_n]^T$ được định nghĩa bằng hàm dấu (sign function):

\begin{equation}
\hat{y} = \text{sign}\left(\sum_{i=1}^{n} w_i x_i + b\right) = \text{sign}(\mathbf{w}^T \mathbf{x} + b)
\end{equation}

trong đó $\mathbf{w}$ là vector trọng số (weights) và $b$ là hệ số độ lệch (bias). Perceptron có thuật toán cập nhật trọng số vô cùng đơn giản dựa trên sai số giữa dự đoán $\hat{y}$ và nhãn thực tế $y$:

\begin{equation}
\mathbf{w}_{new} = \mathbf{w}_{old} + \eta (y - \hat{y})\mathbf{x}
\end{equation}
với $\eta$ là tốc độ học (learning rate). 

Rosenblatt rất tự tin về phát minh của mình, nhưng vào năm 1969, Marvin Minsky và Seymour Papert đã xuất bản cuốn sách "Perceptrons", trong đó họ chứng minh nghiêm ngặt bằng toán học rằng các Perceptron đơn lớp (single-layer perceptron) không thể học được các hàm phân tách phi tuyến tính cơ bản, điển hình là bài toán logic XOR (Exclusive OR). Cuốn sách này đã gây ảnh hưởng lớn đến mức khiến hướng nghiên cứu mạng nơ-ron gần như đóng băng trong suốt nhiều năm.

Phải đến giữa những năm 1980, sự quan tâm đến mạng nơ-ron mới bắt đầu hồi sinh nhờ việc phổ biến hóa thuật toán Lan truyền ngược (Backpropagation). Mặc dù thuật toán này đã được phát minh độc lập bởi nhiều nhà nghiên cứu từ những năm 1960 (như Paul Werbos), nhưng nó chỉ thực sự trở nên nổi tiếng thông qua bài báo đột phá năm 1986 của David Rumelhart, Geoffrey Hinton và Ronald Williams. Lan truyền ngược cung cấp một phương pháp hiệu quả để tính toán gradient (đạo hàm riêng) của hàm mất mát theo từng trọng số trong một mạng nơ-ron đa lớp (Multi-layer Perceptron - MLP), từ đó cho phép sử dụng thuật toán Gradient Descent để tối ưu hóa trọng số:

Giả sử hàm mất mát là $E$, theo quy tắc chuỗi (chain rule) trong giải tích, đạo hàm của $E$ theo trọng số $w_{ij}$ kết nối nơ-ron $j$ ở lớp ẩn với nơ-ron $i$ ở lớp sau là:
\begin{equation}
\frac{\partial E}{\partial w_{ij}} = \frac{\partial E}{\partial o_i} \cdot \frac{\partial o_i}{\partial net_i} \cdot \frac{\partial net_i}{\partial w_{ij}}
\end{equation}
trong đó $net_i$ là tổng đầu vào của nơ-ron $i$, và $o_i = f(net_i)$ là đầu ra sau hàm kích hoạt $f$. Phát hiện này chứng minh rằng mạng nhiều lớp hoàn toàn có khả năng học các hàm phi tuyến tính cực kỳ phức tạp (như XOR), bác bỏ lời chỉ trích của Minsky trước đây.

Năm 1989, Yann LeCun cùng các cộng sự tại AT\&T Bell Labs đã áp dụng thành công thuật toán Backpropagation vào Mạng Nơ-ron Tích chập (Convolutional Neural Networks - CNNs) để nhận dạng chữ số viết tay trên các tấm séc ngân hàng, tạo ra hệ thống LeNet.

Bất chấp những thành công này, trong thập niên 1990 và đầu những năm 2000, mạng nơ-ron lại một lần nữa bị lu mờ trước sự bùng nổ của các thuật toán Học máy (Machine Learning) truyền thống dựa trên thống kê, nổi bật là Support Vector Machines (SVM), Random Forests, và Hidden Markov Models (HMM). Các thuật toán SVM có nền tảng lý thuyết tối ưu hóa toán học vô cùng vững chắc và hoạt động cực kỳ hiệu quả trên các tập dữ liệu nhỏ và trung bình mà không cần tinh chỉnh quá nhiều siêu tham số (hyperparameters). Trong thời kỳ này, trí tuệ nhân tạo trở nên khoa học hơn, có cấu trúc tốt hơn, và dần dần được tích hợp vào các hệ thống khai phá dữ liệu (data mining) và các ứng dụng thương mại một cách thầm lặng.

\newpage

\subsection{Sự bùng nổ của Học sâu (2011 - nay)}

Bước ngoặt thực sự của kỷ nguyên trí tuệ nhân tạo hiện đại diễn ra vào đầu những năm 2010. Mặc dù các cấu trúc mạng nơ-ron đa lớp và thuật toán tối ưu hóa đã có từ nhiều thập kỷ trước, chúng thường gặp phải bài toán vanishing gradient (gradient biến mất) khiến mạng sâu không thể học được, hoặc overfitting do không có đủ dữ liệu. Ba yếu tố mang tính cách mạng đã cùng hội tụ trong giai đoạn này để tạo ra sự bùng nổ của Học sâu (Deep Learning):

\textbf{1. Sự bùng nổ của Dữ liệu lớn (Big Data):} Sự phát triển của Internet, mạng xã hội, và các thiết bị di động đã tạo ra một lượng dữ liệu khổng lồ (hình ảnh, văn bản, âm thanh). Việc có các tập dữ liệu được gán nhãn lớn là điều kiện tiên quyết cho Deep Learning. Điển hình nhất là tập dữ liệu ImageNet do nhóm của giáo sư Fei-Fei Li xây dựng, chứa hàng triệu hình ảnh được gán nhãn thủ công cho hơn một ngàn danh mục đối tượng.

\textbf{2. Sức mạnh phần cứng (GPU):} Các mạng nơ-ron sâu yêu cầu hàng tỷ phép tính ma trận dấu phẩy động. Các Bộ vi xử lý Đồ họa (GPU) do NVIDIA sản xuất, ban đầu được thiết kế để xử lý đồ họa trò chơi điện tử, với kiến trúc xử lý song song hàng nghìn lõi (CUDA cores), tỏ ra hoàn toàn lý tưởng cho việc huấn luyện mạng nơ-ron. Một mô hình từng mất hàng tháng để huấn luyện trên CPU nay có thể hoàn thành trong vài ngày trên GPU.

\textbf{3. Cải tiến Thuật toán và Hàm kích hoạt:} Sự thay thế các hàm kích hoạt truyền thống như Sigmoid hay Tanh bằng hàm ReLU (Rectified Linear Unit), $f(x) = \max(0, x)$, đã giúp giải quyết triệt để vấn đề đạo hàm biến mất. Cùng với đó là các kỹ thuật như Dropout (loại bỏ ngẫu nhiên nơ-ron để chống overfitting) và Batch Normalization.

Sự kiện làm "chấn động" cộng đồng học máy thế giới là cuộc thi phân loại hình ảnh ImageNet (ILSVRC) năm 2012. Nhóm nghiên cứu của giáo sư Geoffrey Hinton tại Đại học Toronto, bao gồm Alex Krizhevsky và Ilya Sutskever, đã giới thiệu AlexNet, một mạng CNN sâu chạy trên GPU. AlexNet đã nghiền nát tất cả các phương pháp thị giác máy tính truyền thống bằng cách giảm tỷ lệ lỗi xuống chỉ còn 15.3\% (so với mức 26.2\% của đối thủ đứng thứ hai). Cột mốc này chứng minh sức mạnh áp đảo của Học sâu trong việc tự động trích xuất các đặc trưng phân cấp từ dữ liệu thô, mở ra kỷ nguyên phục hưng cho AI.

Sau AlexNet, Học sâu tiếp tục càn quét mọi lĩnh vực. Năm 2016, hệ thống AlphaGo do DeepMind phát triển đã đánh bại Lee Sedol, một trong những kỳ thủ cờ vây vĩ đại nhất thế giới, bằng cách kết hợp Học sâu với Học củng cố (Deep Reinforcement Learning). Đây là một cột mốc lịch sử vì cờ vây được cho là trò chơi có không gian trạng thái vô tận (lớn hơn số lượng nguyên tử trong vũ trụ quan sát được) và không thể giải bằng phương pháp tìm kiếm vét cạn bằng máy tính.

\textbf{Cuộc cách mạng Transformer trong Xử lý ngôn ngữ tự nhiên:}
Năm 2017, các nhà nghiên cứu tại Google đã xuất bản bài báo định hình lại toàn bộ lĩnh vực NLP: "Attention Is All You Need". Bài báo giới thiệu kiến trúc \textbf{Transformer}, hoàn toàn loại bỏ các cấu trúc mạng tuần hoàn (RNN/LSTM) trước đó vốn nổi tiếng vì khó song song hóa và hay quên ngữ cảnh xa. Cốt lõi của kiến trúc Transformer là cơ chế \textit{Self-Attention} (tự chú ý), cho phép mô hình tính toán sự tương quan giữa các từ trong câu một cách đồng thời, bất kể vị trí của chúng.

Công thức toán học của cơ chế Scaled Dot-Product Attention được biểu diễn rất thanh lịch:

\begin{equation}
\text{Attention}(Q, K, V) = \text{softmax}\left(\frac{QK^T}{\sqrt{d_k}}\right)V
\end{equation}

trong đó $Q$ (Query), $K$ (Key), và $V$ (Value) là các ma trận đặc trưng được biến đổi từ chuỗi đầu vào, và $d_k$ là kích thước của các vector key, đóng vai trò như một hệ số tỷ lệ để tránh việc hàm softmax bị bão hòa. Cơ chế Multi-Head Attention sau đó kết hợp nhiều không gian biểu diễn khác nhau, cho phép mô hình thu thập thông tin vô cùng phong phú về ngữ pháp và ngữ nghĩa.

Dựa trên kiến trúc Transformer, OpenAI đã phát triển dòng mô hình GPT (Generative Pre-trained Transformer), bao gồm mô hình ngôn ngữ lớn (LLMs). Đồng thời, Google phát triển BERT. Phương pháp tiếp cận "Học tự giám sát" (Self-supervised learning) cho phép các mô hình này học hỏi ngữ nghĩa từ hàng tỷ tài liệu trên Internet bằng cách dự đoán các từ bị che khuất hoặc dự đoán từ tiếp theo mà không cần dữ liệu được gán nhãn của con người.

\newpage

\subsection{Tình trạng hiện tại và Tương lai của Trí tuệ nhân tạo}

Chúng ta đang sống trong kỷ nguyên của Trí tuệ nhân tạo Tạo sinh (Generative AI), với sự ra đời của các Mô hình Ngôn ngữ Lớn (LLMs) như GPT-4, Gemini (Google), Claude (Anthropic), và LLaMA (Meta). Sự xuất hiện của ChatGPT vào cuối năm 2022 được ví như "Khoảnh khắc iPhone" của AI. Lần đầu tiên, một hệ thống AI sở hữu khả năng suy luận, lập trình, viết luận văn sáng tạo, phân tích dữ liệu, và đối thoại lưu loát bằng ngôn ngữ tự nhiên đã được đưa đến tận tay hàng trăm triệu người dùng trên toàn cầu. Khả năng "Zero-shot" và "Few-shot learning" - giải quyết các bài toán mới mà không cần phải tinh chỉnh (fine-tuning) mô hình - đã cách mạng hóa cách chúng ta tương tác với máy tính.

Không chỉ dừng lại ở văn bản, AI hiện đại có khả năng đa phương thức (multimodal). Các hệ thống như Midjourney, DALL-E 3 hay Stable Diffusion có thể tạo ra các tác phẩm nghệ thuật, nhiếp ảnh có độ phân giải siêu cao và độ chân thực cực lớn chỉ từ một câu mô tả (prompt). Trong lĩnh vực y sinh học, hệ thống AlphaFold của DeepMind đã giải quyết thành công bài toán dự đoán cấu trúc protein 3D, mở ra kỷ nguyên mới cho ngành phát hiện và điều chế thuốc (drug discovery).

Mặc dù đạt được những thành tựu phi thường, Trí tuệ nhân tạo hiện tại không phải là không có giới hạn và đối mặt với nhiều thách thức sâu sắc:

\textbf{1. Ảo giác (Hallucination):} Các mô hình ngôn ngữ lớn bản chất là những bộ dự đoán xác suất vô cùng phức tạp, không phải là cơ sở dữ liệu tri thức tĩnh. Chúng có thể tự tin sinh ra các thông tin hoàn toàn sai lệch hoặc bịa đặt các nguồn tài liệu khoa học không tồn tại.

\textbf{2. Khả năng diễn giải (Explainable AI - XAI):} Các mạng nơ-ron sâu với hàng trăm tỷ tham số hoạt động như những "hộp đen" (black box). Các nhà nghiên cứu thường không thể hiểu được cơ chế tại sao một mô hình đưa ra một quyết định hoặc một đầu ra cụ thể. Điều này là một rào cản chí mạng khi áp dụng AI trong các lĩnh vực có tính rủi ro cao như y tế, luật pháp, và xe tự lái, nơi trách nhiệm giải trình (accountability) là bắt buộc.

\textbf{3. Thiên kiến và Đạo đức (Bias and Ethics):} Vì AI được huấn luyện trên dữ liệu lấy từ thế giới thực (vốn luôn chứa đựng những định kiến về chủng tộc, giới tính, tôn giáo), các mô hình thường phản ánh và thậm chí khuếch đại các thiên kiến (bias) này, dẫn đến những phân biệt đối xử tự động trong tuyển dụng hoặc xét duyệt hồ sơ tín dụng. Các vấn đề bản quyền dữ liệu liên quan đến nội dung sáng tạo cũng đang được tranh luận pháp lý rất căng thẳng trên toàn cầu.

\textbf{4. Chi phí môi trường và Phần cứng:} Việc huấn luyện các LLM lớn tốn kém hàng chục triệu đô la và tiêu thụ một lượng điện năng khổng lồ, đi kèm với lượng khí thải carbon cao. Việc tối ưu hóa (quantization, pruning) và phát triển phần cứng xanh, hiệu quả hơn là một hướng nghiên cứu cực kỳ cấp bách.

Tương lai của AI hứa hẹn hướng tới những mô hình suy luận tốt hơn, có thể mô phỏng một phần khả năng lên kế hoạch dài hạn (Agentic AI) và khả năng suy luận logic dựa trên toán học. Cuộc chạy đua công nghệ đang diễn ra gay gắt để tìm kiếm "Trí tuệ nhân tạo tổng quát" (Artificial General Intelligence - AGI) - loại AI có khả năng hiểu, học hỏi và áp dụng kiến thức trong mọi lĩnh vực nhận thức ở mức ngang bằng hoặc vượt trội so với trí tuệ con người. Liệu AGI là một ước mơ kiến tạo nên một kỷ nguyên thịnh vượng mới cho nhân loại hay là một rủi ro hiện sinh (existential risk), như nhận định của nhiều triết gia và nhà khoa học hàng đầu, vẫn là câu hỏi lớn nhất của thế kỷ 21.

\vspace{1cm}
\textbf{Tổng kết:}
Trí tuệ nhân tạo đã đi một chặng đường dài từ những phỏng đoán lý thuyết của Alan Turing và các thuật toán logic của Hội nghị Dartmouth, vượt qua những "mùa đông" khó khăn về kỹ thuật và sự hoài nghi của xã hội, để vươn tới cuộc bùng nổ của sức mạnh tính toán hiện đại và Học sâu. Từ một lĩnh vực nghiên cứu học thuật biệt lập, AI đã trở thành một cơ sở hạ tầng nền tảng, một động lực chính đang tái cấu trúc nền kinh tế, khoa học và xã hội loài người. Chương tiếp theo sẽ đi sâu vào các cấu trúc thuật toán nền tảng tạo nên sức mạnh cho các hệ thống Deep Learning đương đại.
